\songtitle{Tri-Spine País}

\tag*{2026}\tag{carranguera}\tag{thomas-entourage}\tag{spanish-lyrics}
\begin{notes}
Lyrics by SUNO based on the NotebookLM summary of the article \emph{Red Ferroviaria Nacional de Colombia} co-written by Carlos Thompson and Claude.

Part of the \emph{Thomas’ Entourage} narrative universe.
\end{notes}
\begin{style}
Carranguera jingle with fast guacharaca drive, tiple and requinto bouncing in syncopated rural dance feel; verse rides lively storytelling with plucked bass and handclaps, pre-chorus narrows to voice and tiple while the chorus opens with gang-style shout doubles and a catchy rural refrain, Add accordion-like guitar fills, short whistle pickups, and a bright stop-time break before the last chorus, Clean, festive, punchy mix with acoustic sparkle
\end{style}
\begin{lyrics}
\songpart{Verse 1}
De punta a punta va\\
la red a despertar\\
vía vieja, vía nueva\\
se vuelven a encontrar

Por la sierra tan brava\\
la hacen caminar\\
con puentes y andenes\\
la ven renacer

\songpart{Pre-Chorus}
Y en cada estación\\
se siente el empuje\\
carga y pasajero\\
van por el mismo rumbo

\songpart{Chorus}
RFNC, va que va\\
Colombia la va a estrenar\\
RFNC, va que va\\
del puerto pa’ la ciudad\\
Tri-Spine, Tri-Spine\\
la distancia se va a achicar\\
Tri-Spine, Tri-Spine\\
la nación se va a abrazar

\songpart{Verse 2}
Con alianza y visión\\
se arma el porvenir\\
lo público y lo privado\\
se ponen a servir

Aeropuertos, terminales\\
se pueden enlazar\\
menos humo en el camino\\
más vida pa’ sembrar

\songpart{Pre-Chorus}
Cruza el valle, cruza el cerro\\
sin perder el compás\\
cada tren lleva progreso\\
y lo deja al pasar

\songpart{Chorus}
RFNC, va que va\\
Colombia la va a estrenar\\
RFNC, va que va\\
del puerto pa’ la ciudad\\
Tri-Spine, Tri-Spine\\
la distancia se va a achicar\\
Tri-Spine, Tri-Spine\\
la nación se va a abrazar

\songpart{Bridge}
De la costa a la montaña\\
late un mismo andar\\
se juntan pueblos y rutas\\
en un solo cantar

\songpart{Final Chorus}
RFNC, va que va\\
Colombia la va a estrenar\\
RFNC, va que va\\
del puerto pa’ la ciudad\\
Tri-Spine, Tri-Spine\\
la distancia se va a achicar\\
Tri-Spine, Tri-Spine\\
la nación se va a abrazar
\end{lyrics}

